% Suggested replacement for the representation/uniqueness paragraphs in related work.
Shared sum representations underlie Deep Sets~\cite{zaheer2017}. Wagstaff \emph{et al.}\ establish latent-dimension requirements for universal continuous representation and approximation~\cite{wagstaff19,wagstaff22}. Their exact-max argument combines collision analysis with topology; in particular, \cite[Lemma 17]{wagstaff22} treats a fixed cardinality. These results allow particular targets to admit smaller representations. Chen, Chen, and Lu establish representation through multisymmetric polynomial generators and examine readout singularities~\cite{chen2024}.

The closest affine-uniqueness precedent is the restricted Cauchy--Pexider theorem of Charalambides~\cite[Theorem 1.6]{charalambides2013}. For smooth curves without affine subarcs, it identifies affine summands from an additive factorization; its proof regularizes the summands before differentiation. We instead begin with scalar exactness and prove the missing local factorization from a one-way collision implication. Our feature hypothesis also permits $C^1$ curves with affine subintervals. We retain a self-contained proof of this reduction and do not claim affine uniqueness as a new general principle. Scalar monotone-generator uniqueness for quasi-arithmetic means provides a further classical precedent~\cite[Theorems A and D]{matkowski2015}.

% Suggested sentence immediately before Lemma `affine`.
% The affine conclusion is related to restricted Cauchy--Pexider rigidity~\cite[Theorem 1.6 and Sec. V]{charalambides2013}. We prove the version needed here because the collision hypothesis gives only one direction of information preservation and the feature curve is merely $C^1$ and nonaffine.

% Suggested contribution framing.
% First, we derive necessary encoder geometry directly from exact computation of a fixed scalar target. A common fourth input makes encoded triple collisions preserve both moments. Equal-dimensional topology then gives local reverse factorization, allowing an adaptation of the classical restricted Cauchy--Pexider strategy to identify all continuous planar encoders under the stated $C^1$ and nonaffineness assumptions.
